\section{The syntactic recursion category}\label{sec:syntactic}

Let $T$ be a consistent extension of $\PA$ in the language of arithmetic.
Its nonlogical symbols are the constants $0,1$, the binary operations
$+,\cdot$, and the binary relation $<$. Equality is logical. The numeral for $n\in\N$ is the closed term
given by $\bar 0=0$ and $\overline{n+1}=(\bar n+1)$.

In the syntactic recursion category $S_T$ of Montagna and
Di Paola--Montagna~\cite{Montagna,DPM}, we use formal product and
coproduct objects. Let $\mathsf O$ be the set of terms generated
by the constant $\omega$ and the binary operations $\times$ and $+$,
with syntactic equality. Assign a set $|X|$ to each
$X\in\mathsf O$ by
\[
 |\omega|=\N,\qquad
 |X\times Y|=|X|\times|Y|,\qquad
 |X+Y|=|X|\amalg|Y|,
\]
where $A\amalg B=(A\times\{0\})\cup(B\times\{1\})$.
Use the primitive-recursive pairing
$\langle a,b\rangle=2^a(2b+1)-1$ for products and the even and
odd codings for coproducts. These determine recursive bijections
$\gamma_X:|X|\to\N$ by
\[
 \begin{aligned}
 \gamma_\omega&=1_\N,\\
 \gamma_{X\times Y}(x,y)
   &=\langle\gamma_X(x),\gamma_Y(y)\rangle,\\
 \gamma_{X+Y}(x,0)&=2\gamma_X(x),\qquad
 \gamma_{X+Y}(y,1)=2\gamma_Y(y)+1.
 \end{aligned}
\]
The arithmetic formulas below use the single sort of PA.
The numerical pairing, its two projections, and division with
remainder by $2$ are primitive recursive; PA proves their defining
and inverse equations.

Use a deterministic form of Kleene's computation predicate as in
Montagna~\cite[p.~115, Note~1]{Montagna}. Let $E(e,x,s)$ be the primitive-recursive bounded evaluator for program
$e$ on input $x$. At stage $s$ it returns $0$ if no output has been
found and $y+1$ for output $y$.
Write $[a,x]$ for the pairing of the underlying program coding, with
primitive-recursive projections whose inverse equations are provable
in $\PA$. A halting trace has code $t=[s,y]$, where $s$ is its
first successful stage and $y$ its output. The projections are
$\mathsf{st}(t)=s$ and $\mathsf{out}(t)=y$. Thus
\[
 \mathsf{Tr}(e,x,t)\quad\Longleftrightarrow\quad
 E(e,x,\mathsf{st}(t))=\mathsf{out}(t)+1
 \ \land\ \forall r<\mathsf{st}(t)\; E(e,x,r)=0.
\]
The graph of evaluation is
\[
 G(e,x,y)\equiv
 \exists t\,\bigl(\mathsf{Tr}(e,x,t)\land\mathsf{out}(t)=y\bigr),
 \qquad G_e(x,y)=G(\bar e,x,y).
\]
Primitive-recursive symbols abbreviate their defining formulas in
PA, as in \cite[p.~644]{DPM}; the expanded graph formula is $\Sigma_1$.
For $X,Y\in\mathsf O$, let $\mathcal I(X,Y)=\N$.
An index $e\in\mathcal I(X,Y)$ is interpreted using
$\gamma_X$ for inputs and $\gamma_Y$ for outputs.

\begin{lemma}\label{lem:arithmetized-evaluation}
For this computation coding, there are an index $u\in\N$ and
primitive-recursive functions
$\operatorname{s},\operatorname{comp}:\N^2\to\N$
such that PA proves functionality
of $G$ and the universal closures of
\begin{align}
 G_u(\langle e,x\rangle,y)
   &\leftrightarrow G(e,x,y),
   \label{eq:arithmetized-universal}\\
 G(\operatorname{s}(e,a),x,y)
   &\leftrightarrow G(e,\langle a,x\rangle,y),
   \label{eq:arithmetized-specialization}\\
 G(\operatorname{comp}(f,e),x,z)
   &\leftrightarrow
     \exists y\,\bigl(G(e,x,y)\land G(f,y,z)\bigr).
   \label{eq:composition}
\end{align}
Regard $u$ as an element of $\mathcal I(\omega\times\omega,\omega)$.
For each fixed $e\in\N$, there is an index
$k_e\in\N$ for which PA proves
\[
 \forall a\forall p\,
 \bigl(G_{k_e}(a,p)\leftrightarrow
       p=\operatorname{s}(\bar e,a)\bigr).
\]
In particular, PA proves $\forall a\exists!p\,G_{k_e}(a,p)$.
\end{lemma}
\begin{proof}
Use the universal interpreter and formalized parameterization
\cite[p.~108; p.~115, Note~1]{Montagna}. For the program pairing,
the parameter function $s_0$ satisfies
$G(s_0(d,a),x,y)\leftrightarrow G(d,[a,x],y)$ in $\PA$.
On input $\langle e,x\rangle$, let $u$ decode $e,x$ and apply the
interpreter to $[e,x]$. Choose a program $v$ that, on input
$[e,[a,x]]$, computes $\langle a,x\rangle$ and runs $e$ on that input.
The inverse equations for the two pairings give in $\PA$
\[
 G_v([e,[a,x]],y)\longleftrightarrow G(e,\langle a,x\rangle,y).
\]
Two applications of the parameter equation now give
\eqref{eq:arithmetized-specialization} for the primitive-recursive
function $\operatorname{s}(e,a)=s_0(s_0(v,e),a)$.
Choose an index $c$ whose program, on input
$\langle\langle f,e\rangle,x\rangle$, searches successive codes
of pairs $(t,v)$ for which
\[
 \mathsf{Tr}(e,x,t)\land\mathsf{Tr}(f,\mathsf{out}(t),v).
\]
At the first such pair it returns $\mathsf{out}(v)$. Put
$\operatorname{comp}(f,e)=\operatorname{s}(c,\langle f,e\rangle)$,
which is primitive recursive. PA proves that the code of any
witnessing pair bounds the search, and that a successful search
supplies witnesses to both graphs, with intermediate value
$\mathsf{out}(t)$. Two traces for the same input have the same least successful stage
and the same value of $E$ there. Thus $\PA$ proves functionality of
$G$, and all witnessing pairs give the same output, proving~\eqref{eq:composition}.
For $k_e$, use an index for the primitive-recursive function
$a\mapsto\operatorname{s}(e,a)$; its defining equation and
totality are provable in PA.
\end{proof}

The objects of $S_T$ are the terms in $\mathsf O$. Its morphisms
$X\rightharpoonup Y$ are the equivalence classes in
$\mathcal I(X,Y)$ under
\begin{equation}\label{eq:uniform}
 e\equiv_T f
 \quad\Longleftrightarrow\quad
 T\vdash \forall x\forall y\,
       \bigl(G_e(x,y)\leftrightarrow G_f(x,y)\bigr).
\end{equation}
Write $[e]_T$ for the class of $e\in\mathcal I(X,Y)$.

For $e\in\mathcal I(X,Y)$ and $f\in\mathcal I(Y,Z)$, define
$[f]_T[e]_T=[\operatorname{comp}(f,e)]_T$.
Substitution of T-provably equivalent graphs in
\eqref{eq:composition} proves independence of representatives.
Choose an index $i$ for the identity program, so that
$\PA\vdash\forall x\forall y\,(G_i(x,y)\leftrightarrow y=x)$.
Its class in $\mathcal I(X,X)/{\equiv_T}$ is $1_X$.
The identity and associativity equations follow by substituting
the graph of $i$ and reassociating the existential quantifiers
in~\eqref{eq:composition}.

Montagna's objects are the sets generated from $\N$ by
Cartesian products and the displayed disjoint unions
\cite[p.~105]{Montagna}. Let $S_T^{\mathrm{set}}$ denote that
presentation of the unprimed category, using the computation
formulas and the induced numerical codings fixed above
\cite[p.~646]{DPM}. For this set presentation take von Neumann
naturals and Kuratowski ordered pairs. Each $|X|$ is infinite, by
the bijection $\gamma_X$. The element $0=\varnothing$ belongs to
$\N$ and is no ordered pair, so $\N$ is neither a product nor a
tagged union of these sets. The second projection of $A\times B$
is the infinite set $B$, whereas that of $A\amalg B$ is $\{0,1\}$.
The two projections recover the factors of a product, and the two
tag fibres recover the summands of a tagged union. Induction
therefore shows that $|X|=|Y|$ implies $X=Y$. In particular, each
object of the set presentation has a unique coding $\gamma_X$
given by the displayed recursion. There is a functor
\[
 \Pi_T:S_T\longrightarrow S_T^{\mathrm{set}},\qquad
 \Pi_T(X)=|X|,\qquad \Pi_T([e]_T)=[e]_T^{\mathrm{set}}.
\]
Every object of $S_T^{\mathrm{set}}$ is $|X|$ for a term
$X\in\mathsf O$, and the map on each hom-set is a bijection
because both presentations use \eqref{eq:uniform} with the same
numerical codings. Thus $\Pi_T$ is an equivalence. For $T=\PA$,
the target is Montagna's unprimed category $S$ with these choices
of computation formulas and codings.

The pointwise relation used in the primed systems identifies
$e,f\in\mathcal I(X,Y)$ when, for every standard input code $n$,
\[
 T\vdash\forall y\,
   \bigl(G_e(\bar n,y)\leftrightarrow G_f(\bar n,y)\bigr).
\]
Montagna's $S'$ is the case
$T=\PA$ of Di Paola and Montagna's $S'_T$ \cite{Montagna,DPM}.
\emph{Pointwise provable equality and the failure of composition}
\cite[Corollary 2]{Composition} proves that the proposed composition
is not independent of representatives for any consistent c.e.\
extension $T$ of $\PA$. Thus these primed systems do not define
categories with that composition.

For $X\in\mathsf O$ and a $\Sigma_1$ arithmetic formula $C(x)$,
where $x$ is the numerical code variable for $X$, choose
an index $j_C$ for which
\[
 \PA\vdash\forall x\forall y\,
 \bigl(G_{j_C}(x,y)\leftrightarrow(C(x)\land y=x)\bigr).
\]
Such an index searches increasing bounds for a witness to $C(x)$
and returns $x$ when a witness is found
\cite[p.~109, equation~$(*)$]{Montagna}. Write
$[C]_T=[j_C]_T:X\rightharpoonup X$.
For two such formulas,
\begin{equation}\label{eq:domain-order}
 [C]_T\le[D]_T
 \quad\Longleftrightarrow\quad
 T\vdash\forall x\,(C(x)\longrightarrow D(x)).
\end{equation}
Conjunction and disjunction give meets and joins of restriction
idempotents. For $[e]_T:X\rightharpoonup Y$, its restriction
$\overline{[e]_T}$ and range $\ran([e]_T)$ are represented respectively by
\[
 \exists y\,G_e(x,y),\qquad
 R_e(y):=\exists x\,G_e(x,y).
\]
For a restriction idempotent $[C]_T:Y\rightharpoonup Y$, the
equation $[C]_T[e]_T=[e]_T$ holds if and only if
$T$ proves $G_e(x,y)\to C(y)$ universally, equivalently when
$T$ proves $R_e(y)\to C(y)$ universally.
Write $g=[p]_T,h=[q]_T:Y\rightharpoonup Z$.
If $g[e]_T=h[e]_T$, the
composition equation and PA-provable functionality of $G_e$ give
\[
 T\vdash\forall y\forall z\,
 \bigl(R_e(y)\longrightarrow
       (G_p(y,z)\leftrightarrow G_q(y,z))\bigr).
\]
A witness $x$ to $R_e(y)$ can be used in the equality of
the composite graphs, and functionality identifies the intermediate
output with $y$. Hence $g[R_e]_T=h[R_e]_T$ in $S_T$.
The search program $j_{R_e}$ depends only on $e$, so its class
represents the range over every extension of $T$.

For $[e]_T:W\rightharpoonup X$ and $[f]_T:W\rightharpoonup Y$,
a program $h$ runs both computations and pairs their outputs.
For $[p]_T:X\rightharpoonup Z$ and $[q]_T:Y\rightharpoonup Z$,
a program $k$ reads the input tag and runs the indicated program.
Their graphs satisfy the following equivalences in $\PA$:
\[
 \begin{aligned}
 G_h(w,\langle x,y\rangle)&\longleftrightarrow G_e(w,x)\land G_f(w,y),\\
 G_k(2x,z)&\longleftrightarrow G_p(x,z),\\
 G_k(2y+1,z)&\longleftrightarrow G_q(y,z).
 \end{aligned}
\]
Together with the numerical projections and tagged injections,
these give the restriction pairing $[h]_T:W\rightharpoonup X\times Y$
and the copairing $[k]_T:X+Y\rightharpoonup Z$.
Put $\tau_T=[u]_T:\omega\times\omega\rightharpoonup\omega$.
For $f=[e]_T:X\times\omega\rightharpoonup\omega$, put
$r=[k_e]_T:X\to\omega$. Lemma~\ref{lem:arithmetized-evaluation}
proves that $r$ is total and gives a PA-proof of
\[
 G_e(\langle a,x\rangle,y)\leftrightarrow
 \exists p\,
 \bigl(G_{k_e}(a,p)\land G_u(\langle p,x\rangle,y)\bigr)
\]
for all numerical codes $a,x,y$.
The composition and product graph equations therefore give
$f=\tau_T(r\times1_\omega)$ in $S_T$.
Thus $\tau_T$ is a Turing morphism. The class of $i$ in
$S_T(X,\omega)$ represents $\gamma_X$ on the decoded sets.
The same program in $S_T(\omega,X)$ represents $\gamma_X^{-1}$.
Both composites have identity graph by~\eqref{eq:composition}.
The product of any morphism with a nowhere-defined program has
empty graph. Hence $S_T$ is a recursion category.


